\begin{abstract}
    This Bachelor's thesis investigates the potential of Large Language Models (LLMs) for Automated Question Generation (AQG) regarding the adherence to given instructional content and alignment with Bloom's revised Taxonomy. The work is structured around two primary research questions, which will be answered by going through two certain experiments. The methodology integrates both quantitative and qualitative evaluations. Quantitative assessment utilizes cosine similarity measures and LLM-based adherence scoring to assess semantic consistency between generated questions and source content. Qualitative analysis involves expert evaluation using structured rubrics that assess criteria including relevance, clarity, answerability, and especially correctness to the source content in the first experiment. In addition, the second experiment systematically explores the relationship between question formats (Multiple-Choice vs. Open-Ended) and cognitive complexity levels across Bloom's Taxonomy, analyzing how LLMs behave when generating questions with varying instructional information. Results indicate a significant impact of prompt complexity, with common prompts generally outperforming complex ones in terms of e.g.\ content fidelity. Overall, this research demonstrates the promise and challenges of using LLMs for educational question generation and proposes a comprehensive evaluation framework that combines automated metrics with expert assessments to inform future advancements in this emerging field.
\end{abstract}

% #TODO needs to be updated with the final results of the second experiment!

\addcontentsline{toc}{subsection}{Abstract}

\vspace{-1em}

\textcolor{gray}{\hrule}

\vspace{.6em}

\noindent\textbf{Keywords:} Large Language Models; LLMs; Automated Question Generation; AQG; Content Fidelity; Content Adherence; Bloom's Taxonomy; Pedagogical Effectiveness

% \newtodo{Reader should gather the most important information while reading the abstract.}
